\section{Block mollifiers and independent-model mean squares}
\label{sec:model}
\label{model:section}

We seek polynomial multipliers that make the first short piece close
to the positive quantity $A_x$ from \eqref{weights:Ax}, while keeping
the later pieces small. The construction has two separate tasks:
control the multipliers on most integers, and prove mean-square bounds
on the entire independent model. Their combination will allow the
first piece to dominate the weighted-series identity.

Fix a member $F$ of the finite family $\mathcal F$ and a value
$c\in\{-1,1\}$ at the prime $2$, and write $\lambda=\lambda_F$ and $\xi_p=\xi_{F,p}$.
The constants in this section are uniform in these two choices.  We use
the integer characters $z_m$ and the independent variables $Z$ defined
above; in particular, at an odd prime
\[
 \mathbb P(Z(p)=0)=p^{-1},\qquad
 \mathbb P(Z(p)=1)=\mathbb P(Z(p)=-1)=(1-p^{-1})/2,
 \qquad Z(2)=c.
\]
All model norms below are norms on this entire probability space.
We use the parameters and weights of Section~\ref{sec:weights}:
\[
 \ell=\log X,\qquad k_0\le k\le\ell^{3/5},\qquad
 J=\left\lfloor\frac{\log\ell}{10\log2}\right\rfloor,
 \qquad \delta_j=2^{-j},\qquad s=k/\ell.
\]
Here $k$ is an integer.  We always take $X$ sufficiently large that
$J\ge1$ and $0<s\le1$.  The scaled variable $x$ ranges over a fixed
compact subset $\mathcal X$ of $\mathbb R\setminus\{0\}$ containing the
support of the averaging weight.  Constants may depend on this compact
set, the fixed cutoff $f_*$, and $\mathcal F$.

The shift $s$ comes from the weight of the first piece. Replacing the
truncated power in \eqref{weights:W} by its exponential reference gives
the exact integral
\[
 \int_0^\infty e^{-v}
       \exp\bigl(-ku+s\log(C_F|x|v)\bigr)\,dv
       =A_xe^{-ku}.
\]
At $u=\log n/\ell$, the reference weight is therefore $A_xn^{-s}$.
This suggests a multiplier approximating the inverse Euler product at
the real shift $s$. Lemma~\ref{model:first-Fourier-error} and the final
Fourier argument will quantify the replacement in the required norm;
the reference calculation alone does not replace the weighted series.

\subsection{Elementary prime estimates and local factors}

We shall repeatedly use the elementary estimate
$\pi(y)\ll y/\log(2y)$.  One proof starts with
\[
 \sum_{n<p\le2n}\log p\le\log\binom{2n}{n}\le2n\log2.
\]
Summing on dyadic intervals gives $\sum_{p\le y}\log p\ll y$;
separating $p\le\sqrt y$ then proves the asserted estimate for $\pi(y)$.
Partial summation gives the following consequences, with absolute
constants:
\begin{align}
 \sum_{p\le y}\frac1p&\ll1+\log\log(3y),\nonumber\\
 \sum_{a<p\le b}\frac1p&\ll
 1+\log\frac{\log b}{\log a}\quad(2\le a\le b),
 \label{model:prime-mass}\\
 \sum_{p>Y}p^{-1-2a}
 &\ll\int_{2a\log Y}^{\infty}\frac{e^{-v}}v\,dv
 \ll 1+\log^+\frac1{a\log Y}\qquad(0<a\le1,\ Y\ge2).
 \label{model:prime-tail}
\end{align}
Here $\log^+t=\max(0,\log t)$.  If $a\log Y\ge1/2$, we also have
\[
 \int_{2a\log Y}^{\infty}\frac{e^{-v}}v\,dv
 \ll\frac{e^{-2a\log Y}}{a\log Y}.
\]
For clarity, the first bound in \eqref{model:prime-tail} follows by
integrating $t^{-1-2a}$ against $d\pi(t)$, discarding the nonpositive
boundary term at $Y$, and using $1+2a\le3$.
We also need, for $0<a\le1$ and $h\ge a$,
\begin{equation}
 \sum_p\frac{\min(1,h\log p)}p p^{-2a}
 \ll 1+\log(1+h/a).
 \label{model:prime-transition}
\end{equation}
Here is a direct verification.  If $h\le1$, primes with
$\log p\le1/h$ contribute at most
\[
 h\sum_{p\le e^{1/h}}\frac{\log p}{p}\ll1.
\]
Between $1/h$ and $1/a$ the reciprocal-prime mass is
$O(1+\log(h/a))$ by \eqref{model:prime-mass}; primes beyond $e^{1/a}$
contribute $O(1)$ by \eqref{model:prime-tail}.  Empty ranges are omitted.
If $h>1$, discard the minimum and use the same estimates to obtain
$O(1+\log(1/a))$, which is bounded by the right side of
\eqref{model:prime-transition}.

For any real character value $V(p)\in\{-1,0,1\}$ define
\begin{equation}
 D_p(b,V)=1-\lambda(p)V(p)p^{-1/2-b}
                +\xi_p V(p)^2p^{-1-2b},\qquad \Re b\ge0.
 \label{model:inverse-local}
\end{equation}
The local parameter bound gives a factorization
$D_p(b,V)=\prod_{\nu=1}^2(1-\alpha_{p,\nu}V(p)p^{-1/2-b})$,
where $|\alpha_{p,\nu}|\le1$ and zero parameters are permitted.
Consequently
\begin{equation}
 (1-p^{-1/2})^2\le |D_p(b,V)|\le(1+p^{-1/2})^2.
 \label{model:local-bounds}
\end{equation}
For real $b\ge0$, the value of $D_p(b,V)$ is positive: at a good
prime its two factors are complex conjugates (or both positive real),
and at a bad prime its nontrivial factor is positive.  These facts apply
equally to $V=Z$ and $V=z_m$.

\subsection{Construction, length, and the simultaneous product event}

Radziwi\l\l{} and Soundararajan use disjoint prime blocks, truncated
exponential polynomials, and relative truncation estimates in their
study of central values of quadratic twists
\cite[Sections~3, 8, and~9]{RadziwillSoundararajan2015}.
Here we truncate products of inverse local Euler factors by total
degree and estimate shifted norms over the full integer model defined
in Section~\ref{sec:setup}.

Put
\[
 Y_j=X^{\delta_j^2/\sqrt k},\qquad
 P_j^{\mathrm{eu}}(V)=\prod_{p\le Y_j}D_p(s,V)^{-1}
 \quad(0\le j\le J).
\]
Since $s\log Y_0=\sqrt k$, the prime-tail estimate
\eqref{model:prime-tail} makes $\sum_{p>Y_0}p^{-1-2s}$ exponentially
small. This sum will control the residual Euler product in mean square.
For the later pieces, the available
multiplier length decreases as their lengths approach $X$.
The smaller cutoffs $Y_j$ will keep the truncated products within
the comparison allowance $X^{\eta\delta_j/4}$.
In the estimates for these later pieces, the weight factor $\delta_j^k$
will absorb fixed polynomial losses in $k$ and $2^j$.
In the stated range $\log Y_J\ge\ell^{1/2}$, so every $Y_j$ exceeds
$2$ for a single sufficiently large $X$ threshold.  For $i\ge0$ let
\[
 I_{ji}=\{p:Y_j^{2^{-i-1}}<p\le Y_j^{2^{-i}}\},\qquad
 m_{ji}=\lceil k^{1/4}+j+i\rceil.
\]
We omit empty blocks.  They form a disjoint partition of the primes at
most $Y_j$, and \eqref{model:prime-mass} gives
\begin{equation}
 \sum_{p\in I_{ji}}p^{-1}\le C
 \label{model:block-mass}
\end{equation}
uniformly in $j,i,k,X$.  Indeed, when the lower endpoint is at least
$2$, the ratio of the logarithms of the endpoints is $2$; otherwise a
nonempty block has upper endpoint less than $4$.
Set
\begin{align}
 K_{ji}(V,z)&=\prod_{p\in I_{ji}}
  \left(1-\lambda(p)V(p)p^{-1/2-s}z
          +\xi_pV(p)^2p^{-1-2s}z^2\right),
 \nonumber\\
 T_{ji}(V)&=\sum_{b=0}^{m_{ji}}[z^b]K_{ji}(V,z),
 &H_j(V)&=\prod_i T_{ji}(V).
 \label{model:mollifier-definition}
\end{align}
The coefficient extraction in this definition is followed by evaluation
at $z=1$.

\begin{lemma}[Length and coefficients]\label{model:mollifier-length}
The polynomial $H_j$ has an expansion
\[
 H_j(V)=\sum_{r\le R_j}h_j(r)V(r),\qquad
 |h_j(r)|\le\tau(r)r^{-1/2},\qquad
 \sum_r|h_j(r)|\ll R_j^2,
\]
where, for a fixed constant $C_{\rm len}>0$, one may take
\begin{equation}
 R_j=\exp\left(C_{\rm len}(k^{1/4}+j+1)
                     \ell\delta_j^2/\sqrt k\right).
 \label{model:length-bound}
\end{equation}
For every fixed $\eta>0$, a sufficiently large fixed $k_0$ ensures
$R_j\le X^{\eta\delta_j/4}$ for all $0\le j\le J$ and all $k$ in
the stated range.  Thus $H_j$ satisfies the multiplier hypotheses of
Proposition~\ref{analytic:comparison}, with $\Delta=\delta_j/4$.
\end{lemma}

\begin{proof}
In a block, a chosen monomial has prime exponents $e_p\in\{0,1,2\}$
whose sum is at most $m_{ji}$.  Its integer index is therefore at most
$\exp(m_{ji}2^{-i}\log Y_j)$.  Since
\[
 \sum_{i\ge0}m_{ji}2^{-i}
 \le\sum_{i\ge0}(k^{1/4}+j+i+1)2^{-i}
 =2k^{1/4}+2j+4,
\]
we obtain \eqref{model:length-bound}.
Unique factorization shows that each index occurs in exactly one
prime-exponent pattern.  The truncations only delete patterns.  A
surviving coefficient is a product of local coefficients $1$,
$-\lambda(p)p^{-1/2-s}$, or $\xi_pp^{-1-2s}$, and hence is bounded by
$\tau(r)r^{-1/2-s}\le\tau(r)r^{-1/2}$.  For example,
\[
 \sum_{r\le R}\tau(r)r^{-1/2}
 =\sum_{ab\le R}(ab)^{-1/2}
 \le2\sqrt R\sum_{a\le R}\frac1a
 \ll\sqrt R\log(2R)\ll R^2\qquad(R\ge1).
\]
Finally,
\[
 \frac{\log R_j}{\ell\delta_j/4}
 \le4C_{\rm len}\bigl(k^{-1/4}+(j+1)2^{-j}k^{-1/2}\bigr)
 \le8C_{\rm len}k^{-1/4}.
\]
Here $(j+1)2^{-j}\le1$. Thus the fixed choice
$k_0\ge(8C_{\rm len}/\eta)^4$ suffices simultaneously for all $j$.
The coefficient estimate and this length bound verify both multiplier
conditions of Proposition~\ref{analytic:comparison}.
\end{proof}

\begin{lemma}[Prime-sum moments for the integer model]
\label{model:integer-moments}
Let $\mathcal I$ be any set of primes at most $Y_0$, and put
\[
 B_{\mathcal I}(V)=\sum_{p\in\mathcal I}
                 \lambda(p)V(p)p^{-1/2-s}.
\]
If $k_0$ is a sufficiently large absolute constant, then
\begin{equation}
 \frac1X\sum_{X\le|m|<2X}|B_{\mathcal I}(z_m)|^{64}
 \ll\left(1+\sum_{p\in\mathcal I}\frac1p\right)^{32}
 \label{model:64th-moment}
\end{equation}
uniformly in $\mathcal I,k,X$.
\end{lemma}

\begin{proof}
First omit the deterministic prime $2$ and write
$a_p=\lambda(p)p^{-1/2-s}$.  On expanding the model moment, any odd
prime that occurs just once has zero expectation.  Group the remaining
terms by their partition of the $64$ positions into $q\le32$ parts of
sizes at least two.  Since $|a_p|\le2$, for every part of size $r\ge2$
we have $|a_p|^r\le2^{r-2}|a_p|^2$.  Summing the distinct-prime choices
and then discarding distinctness bounds the moment by
\[
 C\sum_{q=1}^{32}\left(\sum_{p\in\mathcal I}|a_p|^2\right)^q
 \ll\left(1+\sum_{p\in\mathcal I}\frac1p\right)^{32}.
\]
Restoring the prime $2$ changes only the constant, by
$|u+v|^{64}\le2^{63}(|u|^{64}+|v|^{64})$ and $|a_2c|\le\sqrt2$.

For completeness the passage to integers needs no estimate for primes
in progressions.  For each ordered $64$-tuple of primes, the product
of its character values is periodic modulo the product of its odd
primes, which is at most $Y_0^{64}$.  Its absolute value is at most one.
The arithmetic mean over the two integer intervals $X\le|m|<2X$ differs from its
complete-residue mean by $O(Y_0^{64}/X)$.  The complete-residue mean is
exactly the independent-model mean, by the Chinese remainder theorem
and the counts of $0,1,-1$ for the Legendre symbol.  There are at most
$Y_0^{64}$ tuples and each coefficient has absolute value at most
$2^{64}$.  Thus the discrepancy between that arithmetic mean and the
model moment is
\[
 O(Y_0^{128}/X)=O(X^{-1+128/\sqrt k})=O(X^{-1/2})
\]
after requiring $\sqrt{k_0}\ge256$.  Since the number of integers in
these intervals is $O(X)$, this proves \eqref{model:64th-moment} with
its normalization by $X$.  If endpoints are not integers, the same
period-counting argument changes only its absolute constant.
\end{proof}

\begin{proposition}[Simultaneous product event]\label{model:good-event}
Except for $O(Xk^{-2})$ integers in $X\le|m|<2X$, the following
inequalities hold simultaneously:
\begin{align}
 \frac12\le H_j(z_m)P_j^{\mathrm{eu}}(z_m)&\le2
                         &&(0\le j\le J),
 \label{model:good-mollifier}\\
 \frac{P_0^{\mathrm{eu}}(z_m)}{P_j^{\mathrm{eu}}(z_m)}
   &\ge2^{-jk/8} &&(1\le j\le J),
 \label{model:good-ratio}\\
 P_0^{\mathrm{eu}}(z_m)&\ge2^{-Jk/8}.
 \label{model:good-full-product}
\end{align}
The implied constant and the lower threshold $k_0$ do not depend on
$k,j,X$.
\end{proposition}

\begin{proof}
For a block $I=I_{ji}$ write $B=B_I(z_m)$.  The elementary expansion
$\log(1+w)=w+O(|w|^2)$ at all sufficiently large primes, together with
\eqref{model:block-mass}, gives
\begin{equation}
 \sup_{|z|=2}|K_{ji}(z_m,z)|\le e^{2|B|+C},
 \qquad K_{ji}(z_m,1)\ge e^{-|B|-C}.
 \label{model:pointwise-block}
\end{equation}
The finitely many smaller primes are included using
\eqref{model:local-bounds} at $z=1$ and the direct upper bound
$\prod_{\nu}(1+2p^{-1/2})$ on $|z|=2$; they change $C$ by an absolute
amount.  Thus no nonvanishing assertion on the radius-two circle is
required.  Cauchy's coefficient estimate now gives, with $m=m_{ji}$,
\[
 \frac{|T_{ji}(z_m)-K_{ji}(z_m,1)|}{K_{ji}(z_m,1)}
 \le C2^{-m}e^{3|B|}.
\]
If every block satisfies $|B_{I_{ji}}(z_m)|\le m_{ji}/100$, the last
quantity is at most $Ce^{-c m_{ji}}$ for a fixed $c>0$.  The sum over
$i$ is $O(e^{-c(k^{1/4}+j)})$.  By enlarging $k_0$, the product over
blocks of the ratios $T_{ji}(z_m)/K_{ji}(z_m,1)$ lies between $1/2$
and $2$.
Since $\prod_iK_{ji}(z_m,1)=(P_j^{\mathrm{eu}}(z_m))^{-1}$,
this proves \eqref{model:good-mollifier} on the specified event.
The number excluded, divided by $X$, is at most
\[
 C\sum_{j=0}^J\sum_{i\ge0}(k^{1/4}+j+i)^{-64}
 \ll (k^{1/4})^{-62}=k^{-31/2}\ll k^{-2},
\]
by Lemma~\ref{model:integer-moments} and
\eqref{model:block-mass}.  The sum may be extended to all $i$ because
only nonempty blocks impose an event.

The real logarithm is available by positivity, and the same local
expansion, including the finitely many small primes separately, gives
for every prime set $\mathcal I$
\begin{equation}
 \log\prod_{p\in\mathcal I}D_p(s,z_m)^{-1}
 =B_{\mathcal I}(z_m)+O\left(\sum_{p\in\mathcal I}p^{-1}\right).
 \label{model:log-product}
\end{equation}
For $Y_j<p\le Y_0$, \eqref{model:prime-mass} bounds this reciprocal
mass by $C(1+j)$, since $\log Y_0/\log Y_j=4^j$.
For $p\le Y_0$ it is at most $C(1+J)$: indeed
$\log\log Y_0\le\log\ell$, and
$\log\ell<10(J+1)\log2$.
Require that the corresponding prime sums have absolute values at
most $(\log2)jk/16$ and $(\log2)Jk/16$, respectively.
Once $k_0$ dominates the fixed drift constants in
\eqref{model:log-product}, these requirements imply
\eqref{model:good-ratio} and \eqref{model:good-full-product}.
Their total exceptional proportion is at most
\[
 C\sum_{j=1}^J\frac{(1+j)^{32}}{(jk)^{64}}
   +C\frac{(1+J)^{32}}{(Jk)^{64}}
 \ll k^{-64}.
\]
Combining the exceptional sets proves the proposition.
\end{proof}

\subsection{Euler-product norms on the full model}

The preceding proposition controls the multipliers on most integers.
We now prove the model norms needed for arithmetic comparison, including
the contribution of every event. We first estimate exact Euler products,
and then replace each complete block by its degree truncation with a
summable relative error.

To avoid any limiting interchange, fix for now a finite real cutoff
$Q\ge Y_0$ and put
\begin{equation}
 \mathcal L_Q(b)=\prod_{p\le Q}D_p(b,Z)^{-1},\qquad \Re b>0.
 \label{model:finite-Euler}
\end{equation}
All bounds in this subsection are independent of $Q$.  At fixed $Q$
the product has the absolutely convergent expansion
\begin{equation}
 \mathcal L_Q(b)=
 \sum_{\substack{n\ge1\\p\mid n\Rightarrow p\le Q}}
       \lambda(n)Z(n)n^{-1/2-b}.
 \label{model:finite-expansion}
\end{equation}
Absolute convergence follows by multiplying the finitely many local
geometric series; their absolute sums are bounded by
$\prod_{p\le Q}(1-p^{-1/2-\Re b})^{-2}$.

\begin{lemma}[Exact shifted products]\label{model:exact-shifts}
Let $0<s,\alpha\le1$, $b=\alpha+it$, and $Y_j\ge2$.  If
$h=\alpha+s+|t|$ and $a=\min(\alpha,s)$, then
\begin{equation}
 \|(P_j^{\mathrm{eu}})^{-1}\mathcal L_Q(b)\|_2
 \ll (1+h/a)^C
       \left(1+\frac1{\alpha\log Y_j}\right)^C.
 \label{model:general-exact-norm}
\end{equation}
At the real point $b=s$, with $j=0$ and the parameters of this section,
\begin{equation}
 \|(P_0^{\mathrm{eu}})^{-1}\mathcal L_Q(s)-1\|_2^2
 \ll \frac{e^{-2\sqrt k}}{\sqrt k}.
 \label{model:exact-distance}
\end{equation}
\end{lemma}

\begin{proof}
We first give the local cancellation with its dependence on $b$.
Write $r=p^{-1/2}$, $u=p^{-s}$, $v=p^{-b}$, and
$M=\min(1,h\log p)$.  The mean-value formula for the exponential and
the trivial bound imply
\begin{equation}
 |u-v|\ll M(u+|v|).
 \label{model:local-difference}
\end{equation}
By \eqref{model:local-bounds}, uniformly for every prime,
\[
 D_p(b,Z)^{-1}-1=O(r|v|).
\]
Subtract the two inverse polynomials and use this last estimate to
obtain
\begin{equation}
 \frac{D_p(s,Z)}{D_p(b,Z)}
 =1+\lambda(p)Z(p)r(v-u)
      +O\big(r^2|v-u|(u+|v|)\big).
 \label{model:local-ratio-expansion}
\end{equation}
For odd $p$, the displayed linear term has expectation zero.  Expanding
the squared absolute value, and using $M^2\le M$ and
$(u+|v|)^2\le2(u^2+|v|^2)$, therefore gives
\begin{equation}
 \mathbb E\left|\frac{D_p(s,Z)}{D_p(b,Z)}\right|^2
 \le \exp\left\{\frac Cp\min(1,h\log p)
                    (p^{-2s}+p^{-2\alpha})\right\}.
 \label{model:local-shift-cost}
\end{equation}
The constants in the error in \eqref{model:local-ratio-expansion} are
uniform even at the small primes, because
$(1-2^{-1/2})^2$ is a common positive lower bound for the denominator.
At the deterministic prime $2$ the local ratio has a fixed upper bound
by \eqref{model:local-bounds}.

For odd primes outside the mollifier, the expansion
\[
 D_p(b,Z)^{-1}=1+\lambda(p)Z(p)p^{-1/2-b}
                       +O(p^{-1-2\alpha})
\]
and centering give
\begin{equation}
 \mathbb E|D_p(b,Z)^{-1}|^2
       \le \exp(Cp^{-1-2\alpha}).
 \label{model:unmatched-local-cost}
\end{equation}
Independence multiplies the local second moments exactly.  Sum
\eqref{model:local-shift-cost} by
\eqref{model:prime-transition}, using both $a=s$ and $a=\alpha$
there, and sum \eqref{model:unmatched-local-cost} by
\eqref{model:prime-tail}.  Taking a square root proves
\eqref{model:general-exact-norm}.

For the distance estimate there is exact cancellation below $Y_0$.
Set
\[
 U=(P_0^{\mathrm{eu}})^{-1}\mathcal L_Q(s)
       =\prod_{Y_0<p\le Q}D_p(s,Z)^{-1},\qquad
 \Xi=\sum_{p>Y_0}p^{-1-2s}.
\]
This is a positive real random variable.  The same expansion gives
\[
 \mathbb E D_p(s,Z)^{-1}=1+O(p^{-1-2s}),\qquad
 \mathbb E D_p(s,Z)^{-2}=1+O(p^{-1-2s}).
\]
The first error may have either sign.  Nevertheless, the elementary
product bound $|\prod(1+e_p)-1|\le\exp(\sum|e_p|)-1$ and
independence show that
\begin{equation}
 \mathbb EU=1+O(\Xi),\qquad \mathbb EU^2=1+O(\Xi).
 \label{model:residual-moments}
\end{equation}
Indeed \eqref{model:prime-tail} and $s\log Y_0=\sqrt k$ give
\[
 \Xi\ll\int_{2\sqrt k}^{\infty}e^{-v}\frac{dv}v
       \ll e^{-2\sqrt k}/\sqrt k,
\]
so $\Xi$ is uniformly small after choosing $k_0$.  Finally,
\[
 \|U-1\|_2^2=\mathbb EU^2-2\mathbb EU+1\ll\Xi.
\]
This uses both moment estimates in \eqref{model:residual-moments} and
proves \eqref{model:exact-distance}.
\end{proof}

\begin{lemma}[Relative truncation in mean square]
\label{model:relative-truncation}
For every $b$ with $\Re b>0$, uniformly in $Q\ge Y_0$,
\begin{equation}
 \|(H_j-(P_j^{\mathrm{eu}})^{-1})\mathcal L_Q(b)\|_2
 \le \left(\exp\left(C\sum_i2^{-m_{ji}}\right)-1\right)
       \|(P_j^{\mathrm{eu}})^{-1}\mathcal L_Q(b)\|_2.
 \label{model:relative-error}
\end{equation}
In particular the relative factor is
$O(2^{-k^{1/4}-j})$.
\end{lemma}

\begin{proof}
For a block $I=I_{ji}$ let
$\mathcal L_I(b)=\prod_{p\in I}D_p(b,Z)^{-1}$.  We claim there are
fixed positive constants $c_*,C_*$ such that
\begin{equation}
 \|K_{ji}(Z,1)\mathcal L_I(b)\|_2\ge c_*,\qquad
 \sup_{|z|=2}\|K_{ji}(Z,z)\mathcal L_I(b)\|_2\le C_*.
 \label{model:block-two-norms}
\end{equation}
For the first assertion, at a large odd prime the exact local ratio is
\[
 \frac{D_p(s,Z)}{D_p(b,Z)}
 =1+\lambda(p)Z(p)p^{-1/2}(p^{-b}-p^{-s})+O(p^{-1}).
\]
The centered linear term shows that its local second moment is
$1+O(p^{-1})$; for all sufficiently large primes it is at least
$1-C/p\ge1/2$.  Their product is bounded below by
$\exp(-2C\sum_{p\in I}1/p)$, which is a fixed positive constant by
\eqref{model:block-mass}.  At each of the finitely many smaller primes,
including $2$, \eqref{model:local-bounds} gives the pointwise lower bound
\[
 \left|\frac{D_p(s,Z)}{D_p(b,Z)}\right|
 \ge\frac{(1-p^{-1/2})^2}{(1+p^{-1/2})^2}>0.
\]
Their finite product is also bounded below.  This proves the first
assertion of \eqref{model:block-two-norms}.

For the second, on $|z|=2$ the local quotient is
\[
 \frac{1-\lambda(p)Z(p)p^{-1/2-s}z
                +\xi_p Z(p)^2p^{-1-2s}z^2}{D_p(b,Z)}
 =1+\lambda(p)Z(p)p^{-1/2}(p^{-b}-zp^{-s})+O(p^{-1})
\]
at large odd primes, uniformly in $z$ and $b$.  Its second moment is
$1+O(p^{-1})$, and therefore its product is bounded above, again by
\eqref{model:block-mass}.  At small primes use the direct numerator
upper bound and \eqref{model:local-bounds}.  The numerator is permitted
to vanish; only an upper bound has been used on this circle.

Apply Cauchy's coefficient formula to the polynomial
$z\mapsto K_{ji}(Z,z)\mathcal L_I(b)$ with values in $L^2$.
For $m=m_{ji}$, Minkowski's inequality and
\eqref{model:block-two-norms} give
\[
 \|(T_{ji}(Z)-K_{ji}(Z,1))\mathcal L_I(b)\|_2
 \le C_*\sum_{h>m}2^{-h}\le C_*2^{-m}.
\]
Dividing by the positive lower bound in
\eqref{model:block-two-norms} makes this a relative error
$\varepsilon_i=C2^{-m_{ji}}$.

To see explicitly how relative errors combine, put
$A_i=K_{ji}(Z,1)\mathcal L_I(b)$ and
$B_i=T_{ji}(Z)\mathcal L_I(b)$.  Then
$\|B_i-A_i\|_2\le\varepsilon_i\|A_i\|_2$ and
$\|B_i\|_2\le(1+\varepsilon_i)\|A_i\|_2$.
Variables belonging to distinct blocks and to the residual primes
$Y_j<p\le Q$ are independent.  In the telescoping identity for
$\prod B_i-\prod A_i$, the norm of each product of factors on distinct
blocks is the product of their norms.  Thus its relative norm, also
including the independent residual product, is at most
\[
 \sum_i\varepsilon_i\prod_{h<i}(1+\varepsilon_h)
   =\prod_i(1+\varepsilon_i)-1
   \le\exp\left(\sum_i\varepsilon_i\right)-1.
\]
This proves \eqref{model:relative-error}, without a factor equal to
the number of blocks.  Finally
$\sum_i2^{-m_{ji}}\le2^{1-k^{1/4}-j}$.
\end{proof}

\begin{proposition}[The shifted norms]\label{model:shifted-norms}
There are fixed $C,c_0>0$ such that, uniformly for every real $\theta$
and every cutoff $Q\ge Y_0$,
\begin{align}
 \|H_0(Z)\mathcal L_Q(s/2+i\theta s)\|_2
       &\ll(1+|\theta|)^C,
 \label{model:half-shift}\\
 \|H_0(Z)\mathcal L_Q(s)-1\|_2
       &\ll e^{-c_0 k^{1/4}},
 \label{model:mollified-distance}\\
 \|H_j(Z)\mathcal L_Q((1+i\theta)/\ell)\|_2
       &\ll (1+|\theta|)^C(k2^j)^C
          \qquad(0\le j\le J).
 \label{model:unit-shift}
\end{align}
\end{proposition}

\begin{proof}
In \eqref{model:general-exact-norm}, first take
$\alpha=s/2$, $t=\theta s$, and $j=0$.
Then $h/\min(\alpha,s)=3+2|\theta|$ and
$\alpha\log Y_0=\sqrt k/2$.  This proves
\eqref{model:half-shift} with the exact multiplier; use
Lemma~\ref{model:relative-truncation} to replace it by $H_0$.
For \eqref{model:unit-shift}, take $\alpha=1/\ell$ and
$t=\theta/\ell$.  Since $k\ge1$,
\[
 h/\min(\alpha,s)=1+k+|\theta|,
 \qquad \alpha\log Y_j=\delta_j^2/\sqrt k.
\]
The bound follows from \eqref{model:general-exact-norm} after increasing
the fixed exponent $C$, and then from relative truncation.
Finally, \eqref{model:exact-distance} bounds the distance from one for
the exact multiplier and also bounds its norm.  The triangle inequality
and \eqref{model:relative-error} give
\[
 \|H_0\mathcal L_Q(s)-1\|_2
 \ll 2^{-k^{1/4}}+e^{-\sqrt k}k^{-1/4}
 \ll e^{-c_0k^{1/4}}.
\]
\end{proof}

\subsection{Fourier representations and the first weight}

We have obtained bounds for the mollified finite Euler products at the
three shifts needed below. Fourier inversion now converts these bounds
into estimates for the actual weights. For the first piece, the main
term is $A_x$. Its estimate combines the power-to-exponential comparison
with the cutoff tail and the exponentially small mollifier error.

We use the Fourier convention
\[
 \widehat g(\theta)=\int_{\mathbb R}g(y)e^{i\theta y}\,dy,
 \qquad
 g(y)=\frac1{2\pi}\int_{\mathbb R}\widehat g(\theta)e^{-i\theta y}
 \,d\theta.
\]
The following elementary bound includes negative real arguments, which
occur inside the gamma integral.

\begin{lemma}[Global approximation of the truncated power]
\label{model:global-power-approximation}
For each fixed nonnegative integer $r$, if $k\ge2r+4$, then for every
$z\in\mathbb R$,
\begin{equation}
 \left|\frac{d^r}{dz^r}
           \left((1-z/k)_+^k-e^{-z}\right)\right|
 \le\frac{C_r}{k}(1+|z|)^2e^{-z}.
 \label{model:global-power-bound}
\end{equation}
\end{lemma}

\begin{proof}
For $z<k$ the derivative of the power is
\[
 (-1)^r\frac{(k)_r}{k^r}(1-z/k)^{k-r},\qquad
 (k)_r=k(k-1)\cdots(k-r+1),
\]
and it is zero for $z\ge k$.  The formula is continuous through $z=k$
in the derivative orders under consideration.  For $z<k$, the logarithm
of its ratio in absolute value to $e^{-z}$ is
\[
 \log((k)_r/k^r)+z+(k-r)\log(1-z/k).
\]
The last two terms have derivative $(r-z)/(k-z)$, so their maximum
over $z<k$ is at $z=r$.  It follows at once that the absolute ratio is
at most $e^r$.  For $|z|\ge\sqrt k$ this bound and
\mbox{$k^{-1}(1+|z|)^2\ge1$} prove \eqref{model:global-power-bound}.
For $|z|<\sqrt k$ we have $|z/k|\le1/2$.  Taylor's formula for
$\log(1-z/k)$ and
$\log((k)_r/k^r)=O_r(k^{-1})$ bound the displayed logarithm by
\[
 O_r\big(k^{-1}(1+|z|)^2\big).
\]
This quantity is bounded in this range; applying
$|e^w-1|\le |w|e^{|w|}$ proves the result there too.  All estimates
hold for negative $z$ as well as positive $z$.
\end{proof}

\begin{lemma}[Fourier error for the first piece]
\label{model:first-Fourier-error}
Define
\[
 A_x=(C_F|x|)^s\Gamma(1+s),\qquad
 E_x(y)=w_{0,x}(y/k)-A_xe^{-y}\quad(y\ge0).
\]
For every required fixed integer $r$, there is an extension $g_x$ of
$e^{y/2}E_x(y)$ from $y\ge0$ to $\mathbb R$ such that
\begin{equation}
 |\widehat g_x(\theta)|\ll_r k^{-1}(1+|\theta|)^{-r}.
 \label{model:first-Fourier-decay}
\end{equation}
The constants are uniform in $x,k,X$ after choosing $k_0$ depending on
$r$.  The quantities $A_x$ are bounded above and below by positive
fixed constants.
\end{lemma}

\begin{proof}
Write $q=C_F|x|$, which lies in a fixed compact subinterval of
$(0,\infty)$.  The exact weight definition gives
\[
 W_{k,x}(y/k)=\int_0^\infty e^{-v}
 \left(1-\frac{y-s\log(qv)}k\right)_+^k\,dv,
\]
whereas
\[
 \int_0^\infty e^{-v}e^{-y+s\log(qv)}\,dv=A_xe^{-y}.
\]
For every fixed $M$, the logarithmic moments
\begin{equation}
 \int_0^\infty e^{-v}(qv)^s
           (1+|s\log(qv)|)^M\,dv\le C_M
 \label{model:gamma-log-moments}
\end{equation}
are uniform for $0<s\le1$: on $0<v\le1$, bound $v^s\le1$ and
$s|\log v|\le|\log v|$; on $v\ge1$, use $v^s\le v$ and the
exponential decay.  Lemma~\ref{model:global-power-approximation} and
differentiation under this dominated integral now imply, for
$0\le h\le r$ and $y\ge-2$,
\begin{equation}
 \left|\frac{d^h}{dy^h}
       \left(W_{k,x}(y/k)-A_xe^{-y}\right)\right|
 \ll_r k^{-1}(1+y^2)e^{-y}.
 \label{model:gamma-power-error}
\end{equation}
In particular the proof uses \eqref{model:global-power-bound} when
$y-s\log(qv)$ is negative and arbitrarily large in absolute value.

The cutoff $F_0(y/k)$ is one for $y\le k/2$, zero for
$y\ge3k/4$, and its $h$th derivative in $y$ is $O_h(k^{-h})$.
Write
\[
 E_x(y)=F_0(y/k)\big(W_{k,x}(y/k)-A_xe^{-y}\big)
                    +(F_0(y/k)-1)A_xe^{-y}
\]
also for $y\ge-2$.  By \eqref{model:gamma-power-error}, after
multiplication by $e^{y/2}$ the first term and its first $r$ derivatives
have $L^1([-2,\infty))$ norm $O_r(k^{-1})$.  The second term is
supported in $y\ge k/2$ and has these norms $O_r(e^{-k/4})$.
Choose a fixed smooth function $\rho$ equal to zero for $y\le-2$ and
one for $y\ge-1$, and extend by
\[
 g_x(y)=\rho(y)e^{y/2}E_x(y),
\]
with zero value to the left of $-2$.  All derivatives through order
$r$ have $L^1(\mathbb R)$ norm $O_r(k^{-1})$.  Integration by parts
$r$ times and the bound without integration by parts give
\eqref{model:first-Fourier-decay}.  Lastly, continuity and positivity
of $q^s\Gamma(1+s)$ on the compact parameter set $q\in[q_-,q_+]$,
$0\le s\le1$, give the two bounds for $A_x$.
\end{proof}

\begin{proposition}[Independent-model estimates for the weight pieces]
\label{model:main-model-estimates}
With the weights and sums $S_{j,Z}(x)$ defined in
Section~\ref{sec:weights}, there is a fixed exponent $C_5$ such
that
\begin{align}
 \|H_0(Z)S_{0,Z}(x)-A_x\|_2&\ll k^{-1},
 \label{model:first-piece-norm}\\
 \|H_j(Z)S_{j,Z}(x)\|_2&\ll(k2^j)^{C_5}\delta_j^k
                    \qquad(1\le j\le J).
 \label{model:later-piece-norm}
\end{align}
These bounds are uniform in $x\in\mathcal X$ and over the entire
range $k_0\le k\le\ell^{3/5}$, with one fixed $k_0$ and one
sufficiently large $X$ threshold.
\end{proposition}

\begin{proof}
Choose a finite cutoff $Q$ exceeding $Y_0$ and the lengths of all the
short polynomials $S_{j,Z}$, $0\le j\le J$.  This is the only
requirement on $Q$; Proposition~\ref{model:shifted-norms} is uniform
in its choice.
For $j\ge1$, set $g_{j,x}(u)=e^u w_{j,x}(u)$.  The support is
contained in
$[1-\delta_j,1-\delta_j/4]\subset[1/2,1]$.
The fixed-order derivative bounds in Lemma~\ref{weights:derivatives}
and integration by parts give, for each required fixed $r$,
\begin{equation}
 |\widehat g_{j,x}(\theta)|
 \ll_r (k/\delta_j)^{C_r}\delta_j^k(1+|\theta|)^{-r}.
 \label{model:later-Fourier-decay}
\end{equation}
Fourier inversion and \eqref{model:finite-expansion} give the exact
identity
\begin{equation}
 S_{j,Z}(x)=\frac1{2\pi}\int_{\mathbb R}
      \widehat g_{j,x}(\theta)
      \mathcal L_Q((1+i\theta)/\ell)\,d\theta.
 \label{model:later-Fourier-representation}
\end{equation}
Indeed every nonzero weight coefficient has index below the chosen
cutoff and hence occurs in \eqref{model:finite-expansion}; all other
indices contribute zero after inversion.  At this fixed finite $Q$,
absolute convergence of the Dirichlet expansion and the integrable
kernel justify interchanging the sum and integral.  Multiply by $H_j$
and use Minkowski's inequality, \eqref{model:unit-shift}, and
\eqref{model:later-Fourier-decay}, taking $r>C+2$.  Their integral
is $O((k/\delta_j)^{C_r}(k2^j)^C\delta_j^k)$, which proves
\eqref{model:later-piece-norm} after fixing $C_5$.

For the first piece, Lemma~\ref{model:first-Fourier-error} and the
same finite-product expansion give
\begin{equation}
 S_{0,Z}(x)-A_x\mathcal L_Q(s)
 =\frac1{2\pi}\int_{\mathbb R}\widehat g_x(\theta)
                    \mathcal L_Q(s/2+i\theta s)\,d\theta.
 \label{model:first-Fourier-representation}
\end{equation}
To check the scaling, at an integer index $n$ put
$y=k\log n/\ell=s\log n\ge0$.  Fourier inversion for $g_x$ then
reads
\[
 w_{0,x}(\log n/\ell)-A_x n^{-s}
 =\frac1{2\pi}\int\widehat g_x(\theta)
                          n^{-s/2-i\theta s}\,d\theta,
\]
which is exactly the coefficient identity needed in
\eqref{model:first-Fourier-representation}.
Use \eqref{model:half-shift} and
\eqref{model:first-Fourier-decay} with $r>C+2$ to bound the norm of
\eqref{model:first-Fourier-representation} after multiplication by
$H_0$ by $O(k^{-1})$.  By \eqref{model:mollified-distance},
\[
 \|H_0S_{0,Z}-A_x\|_2
 \le \|H_0(S_{0,Z}-A_x\mathcal L_Q(s))\|_2
        +A_x\|H_0\mathcal L_Q(s)-1\|_2
 \ll k^{-1}+e^{-c_0k^{1/4}}\ll k^{-1}.
\]
This proves \eqref{model:first-piece-norm}.

All smoothness orders used here are fixed after the absolute exponent
in Proposition~\ref{model:shifted-norms}, and before choosing $k_0$.
The derivative estimates for the weights, the inequality
$s\le\ell^{-2/5}$, and the lower bound
$\delta_J\ge\ell^{-1/10}$ then give a single $X$ threshold for the
whole parameter range.  We have not used the event of
Proposition~\ref{model:good-event} anywhere in these model norm
arguments.  Nor is a passage to an infinite Euler cutoff required:
the coefficient identities hold for every sufficiently large finite
$Q$, and all norm estimates are uniform in that cutoff.
\end{proof}
